% Generated draft: requires mathematical and source review.
\documentclass[11pt]{article}
\usepackage[margin=25mm]{geometry}
\usepackage{fontspec}
\setmainfont{Latin Modern Roman}
\usepackage{amsmath,amssymb,mathtools,unicode-math}
\setmathfont{Latin Modern Math}
\usepackage{enumitem,needspace,xurl,hyperref}
\setlength{\parindent}{0pt}
\setlength{\parskip}{5pt}
\setlength{\emergencystretch}{3em}
\newcommand{\N}{\mathbb{N}}
\newcommand{\Z}{\mathbb{Z}}
\newcommand{\Q}{\mathbb{Q}}
\newcommand{\R}{\mathbb{R}}
\newcommand{\C}{\mathbb{C}}
\newcommand{\F}{\mathbb{F}}
\newcommand{\A}{\mathbb{A}}
\newcommand{\PP}{\mathbb{P}}
\newcommand{\E}{\mathbb{E}}
\newcommand{\eps}{\varepsilon}
\newcommand{\dd}{\,\mathrm{d}}
\begin{document}
\section*{UnsolvedMath ID: 10400184}
\textbf{Source problem code:} catalogue entry
\subsection{Definitions and mathematical statement}
An integral homology \(3\)-sphere is a closed \(3\)-manifold \(M\) with \(H_*(M;\Z)\cong H_*(S^3;\Z)\). In the usual oriented category, let \(\mathcal H\) be their equivalence classes. For the source's finite-type filtration \(F_\bullet\) on the vector space generated by \(\mathcal H\), an invariant \(v\) of degree at most \(d\) is a linear functional vanishing on \(F_{d+1}\). The conjecture is that finite-type invariants separate these manifolds:
\[
 \forall M,N\in\mathcal H,\quad
 M\not\cong N\ \Longrightarrow\
 \exists d<\infty\ \exists v\text{ of finite type of degree at most }d:
 v(M)\ne v(N).
\]
Equivalently, equality of all finite-type invariants implies equivalence of the integral homology spheres. The excerpt does not fix the filtration, coefficient ring, or whether equivalence preserves orientation; rational coefficients and oriented equivalence are the standard contextual reading, not additional printed hypotheses. It refers to Conjecture 11.2 without supplying that statement.

The subsequent supplied subsection gives background rather than an additional separation assertion. A finite-type invariant is \emph{primitive} when
\[
 v(M_1\mathbin{\#}M_2)=v(M_1)+v(M_2)
\]
for all integral homology \(3\)-spheres \(M_1,M_2\). Write \(\mathcal A(\varnothing;R)\) for the Jacobi-diagram module over a coefficient ring \(R\), with no external legs and the usual antisymmetry and IHX relations, and \(\mathcal A(\varnothing;R)_{\mathrm{conn}}\) for the submodule spanned by connected trivalent diagrams. The source states the graded-vector-space identification
\[
 \mathcal A(\varnothing;\Q)
 \cong \operatorname{Sym}_{\Q}
       \bigl(\mathcal A(\varnothing;\Q)_{\mathrm{conn}}\bigr).
\]
No explicit dimension problem is supplied by the intervening subsection heading.
\subsection{Short English statement}
Any two distinct integral homology \(3\)-spheres should be distinguished by some finite-type invariant of finite degree. The supplied primitive-invariant and connected-diagram definitions are accompanying background.
\subsection{Sources}
\begin{itemize}
\item T. Ohtsuki, \emph{Problems on invariants of knots and 3-manifolds} (2002), Conjecture 10.9 and following subsection 10.2.3, supplied PDF page 146. \url{https://msp.org/gtm/2002/04/gtm-2002-04-024s.pdf}. Recorded access 2026-07-29; extraction pdf-font-structure. The source also refers to Conjecture 11.2; its text is absent here.
\item Local UnsolvedMath registry, \texttt{data/unsolvedmath/problems.json}, permanent ID \(10400184\), exact source code \texttt{catalogue entry}; research classification OPEN-TRIAGE and supplied summary classification OPEN.
\end{itemize}
\end{document}
